\chapter{Hardware for Local Inference}
\label{ch:hardware}

Choosing hardware for local models starts with memory capacity and speed.
The machines will change, but understanding these limits helps us judge what
is suitable for a particular task.

\section{Memory, not compute}

To generate each token, the model's weights are read through the processor
once. A thousand tokens means reading the weights a thousand times, so speed is
set by memory bandwidth --- and the model has to fit in memory at all before
speed is a question. A graphics card with 16\,GB cannot run a model that needs
40\,GB at any speed; a machine with 128\,GB of slower shared memory can.

We can estimate the memory needed from the number of parameters and their
precision. Each parameter takes 2 bytes at full precision (FP16), 1 byte at
8-bit, and about 0.5--0.7 bytes at 4-bit (Q4). A 7-billion-parameter (7B) model
therefore needs about 14\,GB, 7\,GB or 4--5\,GB. Add 20--30\% for the context
(the \term{KV cache}) and overheads. A 32\,GB machine comfortably runs a Q4 12B
model with a 32K context; a 70B model wants 48\,GB or more (\Cref{fig:memory}).

\begin{figure}[htbp]
\centering
\fig{ch12-memory-bandwidth}
\caption{Memory, not compute. Every weight is read once per token, so the model
must fit in memory first and speed follows from bandwidth; a 7B model needs
14, 7 or 4--5\,GB depending on precision.}
\label{fig:memory}
\end{figure}

\begin{keypoint}
Capacity first: does the model fit? Bandwidth second: how fast can it be read?
Compute third. Buying for raw processing power is the standard mistake.
\end{keypoint}

\section{Unified memory}

In a conventional PC, the graphics card has its own fixed memory, and anything
it works on is copied there first. \term{Unified memory} puts processor and graphics on one shared pool,
so a 64\,GB laptop can give most of it to a model. Examples include Apple
Silicon, AMD's Ryzen AI Max, Qualcomm Snapdragon X and NVIDIA's desktop
Grace Blackwell. That is why Apple hardware is so common in local AI work: an
architectural advantage, not a brand preference. The trade-off is simple. A separate graphics
card has far more bandwidth (1--3\,TB/s against 200--800\,GB/s) and is the only
sensible choice for training, but it holds only 8--48\,GB and draws 300--700\,W.
A shared-memory machine holds 16--512\,GB and draws 30--120\,W. If the model fits
on the card, the card is faster. The shared-memory advantage is that large models
fit at all (\Cref{fig:unified}).

\begin{figure}[htbp]
\centering
\fig{ch12-unified-memory}
\caption{A separate graphics card against unified memory. The card has more
bandwidth but holds 8--48\,GB; a shared pool is slower per byte but holds
16--512\,GB, so large models fit at all.}
\label{fig:unified}
\end{figure}

\section{Machines, and what they run}

Representative machines and their prices at the time of writing are in
\Cref{app:stateofplay}; they will change within a year. The size guide lasts
longer.

As a guide, text classification needs 1--7B models; a small fine-tuned model
often beats a large general one \parencite{chen2025bionlp}. Summarising and extraction need 7--14B.
Scoped coding needs 12--32B with at least a 32K context, while agent work
across several files needs 30--70B. Below that, expect frustration.
Embedding models are under 1B and run
on an ordinary processor; Whisper for transcription is 1.5B (\Cref{fig:sizes}).

\begin{figure}[htbp]
\centering
\fig{ch12-size-guide}
\caption{A size guide. Parameters on a logarithmic scale, with the tasks each
range serves; names change, but the ranges last longer.}
\label{fig:sizes}
\end{figure}

\section{Quantisation}

\term{Quantisation} stores weights at lower precision. Q4\_K\_M, about four bits
per weight, is the default worth starting from; use Q5 or Q6 if you have room
and the task is sensitive; Q2 and Q3 degrade noticeably. For a fixed amount of
memory, a larger model at Q4 beats a smaller one at Q8
\parencite{dettmers2022fourbit}. These averages hide exceptions. At four bits,
small models can lose much more accuracy than 70B-class ones, and the losses grow
with long inputs and with text in languages other than English
\parencite{lee2025quantisation,mekala2025quantlong}. Test the quantised model on
your own documents before relying on it (\Cref{ch:testing}).

\section{Reasoning effort}

Many recent models write a hidden reasoning trace before answering, and
\term{reasoning effort} sets how long it may run. Locally this setting matters
more than the choice of model. \textcite{willison2026qwen38} found that Qwen 3.8 27B
ships with effort at maximum. Asked for an SVG of a circle, it spent minutes
planning an animated ``geometric study''. A drawing that took 21 minutes with
full reasoning took just over two without. The reasoning also filled the default
8,192-token context before any answer appeared. Start with reasoning low or off
and raise it only when the output is wrong --- in his tests one small tool did
need it. Give the model more context than the runtime's default. Turn on
\term{multi-token prediction} where supported: a light add-on drafts the next
few tokens and the full model confirms them in one pass. With this enabled,
llama.cpp ran the same model about 72 per cent faster than the default
LM Studio build.

\section{Models worth having}

Names date quickly, but the categories are more stable. Start with a coding
model of 12--30B, such as those in the Qwen, Gemma and Kimi families.
Qwen 3.8 27B can drive a coding agent \parencite{willison2026qwen38};
\Cref{app:stateofplay} lists versions current at the time of writing.
Check tool-calling support before relying on a model for agents.

For summarising and extraction, keep a general instruct model of 7--30B.
An embedding model (\code{bge-m3}) and a reranker (\code{bge-reranker}) serve
retrieval, while Whisper handles transcription. A small 1--4B model is
useful for high-volume classification. Dense models read every weight for every token, so
memory bandwidth slows them most. Check each licence: several ``open'' models
restrict use.

\section{Shared deployment}

For a group, one good machine beats several mediocre ones: run vLLM on a
128\,GB server, expose an OpenAI-compatible endpoint on the institutional
network, and point everyone's tools at it.

\begin{lstlisting}[style=shell,caption={Minimal shared serving}]
vllm serve Qwen/Qwen3-30B-A3B-Instruct --host 0.0.0.0 --port 8000 \
    --max-model-len 32768                          # on the server
export OPENAI_BASE_URL=http://gpu-server.local:8000/v1   # on each workstation
export OPENAI_API_KEY=not-used-but-required
\end{lstlisting}

Keep it behind the institutional firewall. An open inference endpoint is free
computing for whoever finds it, and it will be found.
